\section{background}
\label{background}
\subsection{SHAP}

\gls{SHAP} is an explainability method grounded in cooperative game theory. It attributes a model’s output to its input features by computing their marginal contributions to the prediction. Originally developed to ensure fair reward distribution in multiplayer games, SHAP applies this principle to machine learning models by treating features as players in a coalition.

In the context of machine learning, SHAP can be used to understand why a model makes a prediction. It provides local explanations that highlight the most influential features for an individual classification and global insights that reflect the model’s overall decision-making behavior. This enhances transparency and trust, especially in high-stakes domains like cyber defense.

Let $f$ denote the original prediction model, and let $g$ represent the explanation model, which approximates $f$ in a human-understandable way. The explanation model operates on simplified inputs $x'$, which are mapped back to the original input space via a transformation $x = h(x')$. 

A common approach is to define $g$ as a linear function of binary variables:
\begin{equation*}
    g(x') = \phi_{0} + \sum_{i=1}^{M} \phi_{i} x'_{i}
\end{equation*}
where M is the number of simplified input features, $\phi_{i}$ denotes the feature attribution (importance weight) of the i-th feature, and $\phi_{0}$ is a baseline value representing the model's output in the absence of any feature. 

Let $F$ be the set of all features, $S$ is a subset of features excluding feature $i$ and $f_{S}(x_{S})$ denotes the model prediction restricted to the feature subset $S$. Then the SHAP value for the i-th feature is given by:
\begin{equation*}
    \phi_i =
\sum_{S \subseteq F \setminus \{i\}}
\frac{|S|!(|F| - |S| - 1)!}{|F|!}
\left( f_{S \cup \{i\}}(x_{S \cup \{i\}}) - f_S(x_S) \right)
\end{equation*}

Training a model for each combination results in a combinatorial explosion. To address this, the SHAP Explainer uses background data, typically a subset of the training dataset. Instead of retraining the model without a specific feature, the feature is replaced with a sample from the background data that reflects the distribution of the training data. This approach enables the calculation of SHAP values without requiring multiple model training runs.